La dérivation de \(G\) porte cette réciprocité à une condition géométrique beaucoup plus concrète.\par

La mécanique dimensionnelle détermine d’abord quel type d’objet \(G\) peut être. Elle fixe sa droite de grandeur et son homogénéité dimensionnelle. La clôture chronologique sélectionne ensuite sa valeur : la dimension détermine la classe de grandeur et la clôture détermine son coefficient adimensionnel.\par

La sélection procède de la périodicité temporelle de (108) phases.\par

Le cycle \(108\) détermine une fréquence, un rayon de phase et une énergie. À partir de cette énergie apparaît une masse chronologique. Et en construisant la longueur de Compton réduite de cette masse, on obtient exactement le même rayon du cycle :\par

\[
\bar\lambda_C=R_{108}.
\]

La phase et la matière expriment déjà la même longueur avant l’introduction de \(G\).\par

La gravité intervient alors. On construit le rayon gravitationnel réduit\par

\[
r_g=\frac{Gm}{c^2},
\]

et l’on exige que la clôture gravitationnelle identifie le rayon de phase, la longueur de Compton et le rayon gravitationnel comme une seule échelle structurelle :\par

\[
r_g=R_{108}=\bar\lambda_C.
\]

Cette condition sélectionne de manière univoque\par

\[
\theta_{108}
=
\left(\frac{54}{\pi}\right)^2.
\]

Ainsi,\par

\[
R_{108}
=
\bar\lambda_C
=
r_g
=
\ell_P.
\]

La fonction structurelle de \(\theta_{108}\) détermine la signification de son nombre décimal terminal : elle montre ce que fait \(G\).\par

\(G\) apparaît comme le coefficient nécessaire pour faire coïncider les quatre longueurs publiées par ces langages :\par

\begin{itemize}[leftmargin=*,itemsep=0.35em]
\item le langage de la phase produit \(R_{108}\) ;
\item le langage quantique produit \(\bar\lambda_C\) ;
\item le langage gravitationnel produit \(r_g\) ;
\item le langage de Planck produit \(\ell_P\).
\end{itemize}

La constante gravitationnelle est le médiateur qui empêche ces quatre lecteurs de se séparer.\par

Dit de manière plus physique : \(G\) exprime combien la géométrie doit répondre pour qu’un quantum dont la masse provient de l’horloge \(108\) se referme sur sa propre phase.\par

C’est la relation qui rend possible que matière, phase et géométrie soient trois publications compatibles du même événement.\par

Ici, la lecture machienne devient plus forte. L’échelle gravitationnelle locale est sélectionnée par une condition de clôture globale du cycle. Le rayon local d’un quantum acquiert un sens en s’intégrant dans une périodicité complète. L’inertie, la phase et la géométrie se déterminent relationnellement.\par

La portée démontrée fournit un mécanisme mathématique concret pour une lecture machienne : une propriété locale — le couplage gravitationnel — est sélectionnée par la cohérence de l’état global. Les autres formulations historiques du principe de Mach conservent leur domaine propre.\par

Ensuite apparaît la dimension cinq.\par

Le scalaire \(G\) fixe l’intensité gravitationnelle et le porteur spécifie sa structure représentationnelle. HMT construit un espace tensoriel symétrique de trace nulle et de dimension cinq. Ce porteur admet cinq poids hélicoïdaux :\par

\[
+2,\quad +1,\quad 0,\quad -1,\quad -2.
\]

Ce sont les cinq degrés de liberté naturels d’un tenseur symétrique sans trace en trois dimensions avant l’imposition des contraintes dynamiques d’une théorie particulière.\par

Cela clarifie un point souvent confondu. « Cinq composantes » désigne le porteur complet, tandis que les deux hélicités désignent la réduction transverse de masse nulle en relativité générale.\par

Cela signifie que le porteur complet possède cinq coordonnées irréductibles. Ensuite, la dynamique décide quelle partie de ce porteur est physique :\par

\begin{itemize}[leftmargin=*,itemsep=0.35em]
\item dans une réalisation massive de spin deux, les cinq poids peuvent subsister ;
\item dans une réalisation sans masse avec invariance de jauge, la projection transverse sans trace élimine les secteurs \(\pm1\) et \(0\) ;
\item les hélicités \(\pm2\) subsistent.
\end{itemize}

Ainsi, HMT fournit un espace antérieur à la décision dynamique. Elle construit le porteur complet à cinq composantes et contient, par la projection correspondante, sa réduction physique à deux hélicités.\par

La beauté réside dans la coexistence ordonnée de deux affirmations :\par

\begin{itemize}[leftmargin=*,itemsep=0.35em]
\item la géométrie de spin deux possède un porteur à cinq composantes ;
\item le rayonnement gravitationnel sans masse occupe un sous-espace bidimensionnel.
\end{itemize}

Les deux affirmations forment une réduction cohérente et précise.\par

Et cette réduction conserve un autre enseignement central : le scalaire \(G\) fixe l’intensité commune et l’information d’orientation réside dans le porteur tensoriel. L’intensité réside dans la constante, les hélicités résident dans la représentation et la généalogie complète réunit grandeur et porteur.\par

